

\begin{theorem}[Dual $L_2$-gain {\cite[Thm.~9]{shivaknew}}]\label{thm:duall2gain}
	For a PIE system with corresponding PI operators $G\{\mcl{T}, \mcl{A}, \mcl{B}, \mcl{C}, \mcl{D}\}$, the following statements are equivalent.
	\begin{enumerate}
		\item For any $w\in\mbf{L}_{e,[a,b]}^{\mbf{n}_w}$, and $\mbf{x}(0)=0$, $G$  satisfies $\|P_\tau {z}\| \leq \|P_\tau {w}\|$.
		\item For any $\underline{w}\in\mbf{L}_{e,[a,b]}^{\mbf{n}_z}$, and $\underline{\mbf{x}}(0)=0$, $G^\top$ satisfies $\|P_\tau \underline{z}\| \leq \|P_\tau \underline{w}\|$.
	\end{enumerate}
\end{theorem}

\begin{definition}[Primal and Dual Dynamic Scalings]
	Let	$\Psi=\bmat{\Psi_z & \Psi_{zw}\\ \Psi_{wz} & \Psi_w}$ satisfy
	$\Psi,\Psi^{-1}\in\RH$. Let $\hat{\mcl{J}}:=\bmat{0&I\\ -I&0}$.
	The dual filter associated with $\Psi$ is defined by the
	dual operation
	\begin{equation}\label{eqn:dualoperation}
		\mbf{D}(\Psi)
		:=
		\left(\hat{\mcl{J}}^*\Psi^\top\hat{\mcl{J}}\right)^{-1}
		=
		\bmat{\Psi_w^\top & -\Psi_{zw}^\top\\
		      -\Psi_{wz}^\top & \Psi_z^\top}^{-1}.
	\end{equation}
	Note $\mbf{D}(\Psi)\in\RH$. Since $\hat{\mcl{J}}^*=\hat{\mcl{J}}^{-1}$,
	$\hat{\mcl{J}}^*=-\hat{\mcl{J}}$, $(\Psi^\top)^\top=\Psi$, and
	$(\Psi^{-1})^{-1}=\Psi$, this operation is an involution:
	\begin{equation}\label{eqn:dualoperationinvolution}
		\mbf{D}(\mbf{D}(\Psi))=\Psi.
	\end{equation}

	The filter $\Psi$ admits the state-space
	realization
	\begin{equation}\label{eqn:ssprimalscaling:iqc}
		\bmat{\mcl{T}_\Psi \dot{\mbf{x}}_\Psi(t)\\ \tilde{z}(t) \\ \tilde{w}(t)}
		=
		\bmat{\mcl{A}_\Psi & \mcl{B}_\Psi^z & \mcl{B}_\Psi^w \\
			\mcl{C}_\Psi^z & \mcl{D}_\Psi^z & \mcl{D}_\Psi^{zw}\\
			\mcl{C}_\Psi^w & \mcl{D}_\Psi^{wz} & \mcl{D}_\Psi^w}
		\bmat{\mbf{x}_\Psi(t) \\ z(t) \\ w(t)},
	\end{equation}
	where the entries are PI operators of appropriate dimension.
	The dual filter $\mbf{D}(\Psi)$ is equivalently represented
	through the inverse relation $\mbf{D}(\Psi)^{-1}=\hat{\mcl{J}}^*\Psi^\top\hat{\mcl{J}}$, which has
	the state-space realization
	\begin{equation}\label{eqn:ssdualscaling:iqc}
		\bmat{\mcl{T}_{\Psi}^*\dot{\underline{\mbf{x}}}_{\Psi}(t)\\
		\underline{z}(t) \\ \underline{w}(t)}
		=
		\bmat{\mcl{A}_\Psi^* & -\mcl{C}_\Psi^{w*} & \mcl{C}_\Psi^{z*} \\
			-\mcl{B}_\Psi^{w*} & \mcl{D}_\Psi^{w*} & -\mcl{D}_\Psi^{zw*} \\
			\mcl{B}_\Psi^{z*} & -\mcl{D}_\Psi^{wz*} & \mcl{D}_\Psi^{z*}}
		\bmat{\underline{\mbf{x}}_\Psi(t) \\
		\tilde{\underline{z}}(t) \\ \tilde{\underline{w}}(t)}.
	\end{equation}
	The filtered primal system is obtained
	from
	\begin{equation}\label{eqn:primalscaling:iqc}
		\Psi \bmat{z(t) \\ w(t)}
		=
		\bmat{\tilde{z}(t) \\ \tilde{w}(t)}.
	\end{equation}
	The filtered dual system is obtained from
	\begin{equation}\label{eqn:dualscaling:iqc}
		\mbf{D}(\Psi)
		\bmat{\underline{z}(t) \\ \underline{w}(t)}
		=
		\bmat{\tilde{\underline{z}}(t) \\ \tilde{\underline{w}}(t)}.
	\end{equation}
\end{definition}

Adopted from \cite{venkataraman}
\begin{definition}(Potapov--Ginzburg Transform)\label{def:pg}
    Let ${\Psi} : \bmat{z \\ w} \to \bmat{\tilde{z} \\ \tilde{w}}$, where $\Psi = \bmat{\Psi_{z} & \Psi_{zw} \\ \Psi_{wz} & \Psi_{w}}$. 
	Then the Potapov-Ginzburg transform of $\Psi$ is defined as
    \begin{equation}\label{eqn:pgtransform}
        \bmat{\tilde{z} \\ w}= \Psi^\dagger \bmat{\tilde{w} \\ z} \text{, where }\Psi^\dagger = \bmat{\Psi_{zw}\Psi_{w}^{-1} & \Psi_{z} - \Psi_{zw}\Psi_{w}^{-1}\Psi_{wz}\\
                            \Psi_{w}^{-1} & -\Psi_{w}^{-1}\Psi_{wz}}.
    \end{equation}
    Assumed to admit the state-space representation
    \begin{equation}\label{eqn:sspgtransform}
        \bmat{\mcl{T}_\Psi\dot{\mbf{x}}_\Psi(t)\\ \tilde{z}(t) \\ w(t)} = 
        \bmat{\hat{\mcl{A}}_\Psi & \hat{\mcl{B}}_\Psi^{\tilde{w}} & \hat{\mcl{B}}_\Psi^z\\
        \hat{\mcl{C}}_\Psi^z & \hat{\mcl{D}}_\Psi^{z\tilde{w}} & \hat{\mcl{D}}_\Psi^z \\
        \hat{\mcl{C}}_\Psi^w & \hat{\mcl{D}}_\Psi^{w\tilde{w}} & \hat{\mcl{D}}_\Psi^{wz}}\bmat{\mbf{x}_\Psi(t) \\ \tilde{w}(t) \\ z(t)},
    \end{equation}
    where
    \begin{align*}
        \hat{\mcl{A}}_\Psi &:= \mcl{A}_\Psi - \mcl{B}_\Psi^w(\mcl{D}_\Psi^w)^{-1}\mcl{C}_\Psi^w, &
        \hat{\mcl{B}}_\Psi^{\tilde{w}} &:= \mcl{B}_\Psi^w(\mcl{D}_\Psi^w)^{-1}, \\
        \hat{\mcl{B}}_\Psi^{z} &:= \mcl{B}_\Psi^z - \mcl{B}_\Psi^w(\mcl{D}_\Psi^w)^{-1}\mcl{D}_\Psi^{wz}, &
        \hat{\mcl{C}}_\Psi^z &:= \mcl{C}_\Psi^z - \mcl{D}_\Psi^{zw}(\mcl{D}_\Psi^w)^{-1}\mcl{C}_\Psi^w,\\
        \hat{\mcl{D}}_\Psi^z &:= \mcl{D}_\Psi^z - \mcl{D}_\Psi^{zw}(\mcl{D}_\Psi^w)^{-1}\mcl{D}_\Psi^{wz}, &
        \hat{\mcl{D}}_\Psi^{z\tilde{w}} &:= \mcl{D}_\Psi^{zw}(\mcl{D}_\Psi^w)^{-1}, \\
        \hat{\mcl{C}}_\Psi^w &:= -(\mcl{D}_\Psi^w)^{-1}\mcl{C}_\Psi^w, &
        \hat{\mcl{D}}_\Psi^{w\tilde{w}} &:= (\mcl{D}_\Psi^w)^{-1}, \\
        \hat{\mcl{D}}_\Psi^{wz} &:= -(\mcl{D}_\Psi^w)^{-1}\mcl{D}_\Psi^{wz}.
    \end{align*}
\end{definition}


\begin{lemma} (Potapov--Ginzburg Transform Equivalence)\label{lem:pgequivalence}
The \textbf{Scaled Primal System} (respectively, \textbf{Scaled Dual
System}) is equivalent, under an invertible congruence transformation,
to the \textbf{Primal System} (respectively, \textbf{Dual System})
augmented with the Potapov--Ginzburg transform of the \textbf{Primal
Scaling} (respectively, \textbf{Dual Scaling}).
\end{lemma}
\begin{proof}
    We first show the primal derivation explicitly. To show
    sufficiency, first augment the \textbf{Primal System} with the
    Potapov--Ginzburg transform of the \textbf{Primal Scaling}.
    Substitution of $z=\mcl{C}\mbf{x}+\mcl{D}w$ into
    \eqref{eqn:sspgtransform} gives the algebraic loop
    \begin{align*}
        w
            &=
            \Lambda_p\hat{\mcl{D}}_\Psi^{wz}\mcl{C}\mbf{x}
            +\Lambda_p\hat{\mcl{C}}_\Psi^w\mbf{x}_\Psi
            +\Lambda_p\hat{\mcl{D}}_\Psi^{w\tilde{w}}\tilde{w}.
    \end{align*}
    Assume that this loop is well-posed, and define
    \begin{equation*}
        \Lambda_p := (I-\hat{\mcl{D}}_\Psi^{wz}\mcl{D})^{-1},
        \qquad
        \Lambda_z := (I-\mcl{D}\hat{\mcl{D}}_\Psi^{wz})^{-1}.
    \end{equation*}
    Solving for $w$ and substituting into the plant state equation gives
    \begin{align*}
        \mcl{T}\dot{\mbf{x}}
            &= \left(\mcl{A}
                +\mcl{B}\Lambda_p\hat{\mcl{D}}_\Psi^{wz}\mcl{C}\right)\mbf{x}
               +\mcl{B}\Lambda_p\hat{\mcl{C}}_\Psi^w\mbf{x}_\Psi
               +\mcl{B}\Lambda_p\hat{\mcl{D}}_\Psi^{w\tilde{w}}\tilde{w}.
    \end{align*}
    Using $\mcl{D}\Lambda_p=\Lambda_z\mcl{D}$ and
    $I+\mcl{D}\Lambda_p\hat{\mcl{D}}_\Psi^{wz}=\Lambda_z$, as verified
    below in the push-through identities, \eqref{eqn:id1} and \eqref{eqn:id4}, we obtain
    \begin{equation*}
        z
        =
        \Lambda_z\mcl{C}\mbf{x}
        +\Lambda_z\mcl{D}\hat{\mcl{C}}_\Psi^w\mbf{x}_\Psi
        +\Lambda_z\mcl{D}\hat{\mcl{D}}_\Psi^{w\tilde{w}}\tilde{w}.
    \end{equation*}
    Substitution into the equations for $\dot{\mbf{x}}_\Psi$ and
    $\tilde{z}$ yields the augmented realization
    \begin{equation}\label{eqn:potapov}
    \bmat{\mcl{T}\dot{\mbf{x}}(t) \\ \mcl{T}_\Psi\dot{\mbf{x}}_\Psi(t)\\ \tilde{z}(t)} = 
    \bmat{\mcl{A} + \mcl{B}\Lambda_p\hat{\mcl{D}}_\Psi^{wz}\mcl{C} & \mcl{B}\Lambda_p\hat{\mcl{C}}_\Psi^w & \mcl{B}\Lambda_p\hat{\mcl{D}}_\Psi^{w\tilde{w}} \\
    \hat{\mcl{B}}_\Psi^z\Lambda_z\mcl{C} & \hat{\mcl{A}}_\Psi + \hat{\mcl{B}}_\Psi^z\Lambda_z\mcl{D}\hat{\mcl{C}}_\Psi^w & \hat{\mcl{B}}_\Psi^{\tilde{w}} + \hat{\mcl{B}}_\Psi^z\Lambda_z\mcl{D}\hat{\mcl{D}}_\Psi^{w\tilde{w}} \\
    \hat{\mcl{D}}_\Psi^z\Lambda_z\mcl{C} & \hat{\mcl{C}}_\Psi^z + \hat{\mcl{D}}_\Psi^z\Lambda_z\mcl{D}\hat{\mcl{C}}_\Psi^w & \hat{\mcl{D}}_\Psi^{z\tilde{w}} + \hat{\mcl{D}}_\Psi^z\Lambda_z\mcl{D}\hat{\mcl{D}}_\Psi^{w\tilde{w}}}\bmat{\mbf{x}(t) \\ \mbf{x}_\Psi(t) \\ \tilde{w}(t)}
    \end{equation}
Define also
\begin{equation*}
    \Lambda_c := (\mcl{D}_\Psi^{wz} \mcl{D} + \mcl{D}_\Psi^w)^{-1} = 
    (I+(\mcl{D}_\Psi^w)^{-1}\mcl{D}_\Psi^{wz} \mcl{D})^{-1}(\mcl{D}_\Psi^w)^{-1} = 
    (\mcl{D}_\Psi^w)^{-1} (I+\mcl{D}_\Psi^{wz}\mcl{D}(\mcl{D}_\Psi^w)^{-1})^{-1}.
\end{equation*}
The push-through identities
\begin{equation*}
    (I+\mcl{U}\mcl{V})^{-1}\mcl{U}
        = \mcl{U}(I+\mcl{V}\mcl{U})^{-1}, \qquad
    (I+\mcl{U}\mcl{V})^{-1}
        = I-\mcl{U}(I+\mcl{V}\mcl{U})^{-1}\mcl{V},
\end{equation*}
give the following relations
\begin{align}
    \Lambda_p(\mcl{D}_\Psi^w)^{-1}
        &= \left(I+(\mcl{D}_\Psi^w)^{-1}\mcl{D}_\Psi^{wz}\mcl{D}\right)^{-1}
           (\mcl{D}_\Psi^w)^{-1} = \Lambda_c
           \label{eqn:id1}\\[4pt]
    (\mcl{D}_\Psi^w)^{-1}\mcl{D}_\Psi^{wz}\Lambda_z
        &= (\mcl{D}_\Psi^w)^{-1}\mcl{D}_\Psi^{wz}
           \left(I+\mcl{D}(\mcl{D}_\Psi^w)^{-1}\mcl{D}_\Psi^{wz}\right)^{-1}
           \notag\\
        &= (\mcl{D}_\Psi^w)^{-1}
           \left(I+\mcl{D}_\Psi^{wz}\mcl{D}(\mcl{D}_\Psi^w)^{-1}\right)^{-1}
           \mcl{D}_\Psi^{wz}
         = \Lambda_c\mcl{D}_\Psi^{wz}
           \label{eqn:id2}\\[4pt]
    \Lambda_z\mcl{D}(\mcl{D}_\Psi^w)^{-1}
        &= \left(I+\mcl{D}(\mcl{D}_\Psi^w)^{-1}\mcl{D}_\Psi^{wz}\right)^{-1}
           \mcl{D}(\mcl{D}_\Psi^w)^{-1}
           \notag\\
        &= \mcl{D}(\mcl{D}_\Psi^w)^{-1}
           \left(I+\mcl{D}_\Psi^{wz}\mcl{D}(\mcl{D}_\Psi^w)^{-1}\right)^{-1}
         = \mcl{D}\Lambda_c
           \label{eqn:id3}\\[4pt]
    \Lambda_z
        &= \left(I+\mcl{D}(\mcl{D}_\Psi^w)^{-1}\mcl{D}_\Psi^{wz}\right)^{-1}
           \notag\\
        &= I-\mcl{D}(\mcl{D}_\Psi^w)^{-1}
           \left(I+\mcl{D}_\Psi^{wz}\mcl{D}(\mcl{D}_\Psi^w)^{-1}\right)^{-1}
           \mcl{D}_\Psi^{wz}
           \notag\\
        &= I-\mcl{D}\Lambda_c\mcl{D}_\Psi^{wz}
           \label{eqn:id4}\\[4pt]
    (\mcl{D}_\Psi^w)^{-1}
        -\Lambda_c\mcl{D}_\Psi^{wz}\mcl{D}(\mcl{D}_\Psi^w)^{-1}
        &= \left(I-\Lambda_c\mcl{D}_\Psi^{wz}\mcl{D}\right)
           (\mcl{D}_\Psi^w)^{-1}
           \notag\\
        &= \Lambda_c\left(\mcl{D}_\Psi^w
           +\mcl{D}_\Psi^{wz}\mcl{D}
           -\mcl{D}_\Psi^{wz}\mcl{D}\right)(\mcl{D}_\Psi^w)^{-1}
           \notag\\
        &= \Lambda_c.
           \label{eqn:id5}
\end{align}
We now express the blocks of \eqref{eqn:potapov} in terms of
$\Lambda_c$.  For the first row, direct substitution and
\eqref{eqn:id1} give
\begin{align*}
&\bmat{
    \mcl{A}+\mcl{B}\Lambda_p\hat{\mcl{D}}_\Psi^{wz}\mcl{C} &
    \mcl{B}\Lambda_p\hat{\mcl{C}}_\Psi^w &
    \mcl{B}\Lambda_p\hat{\mcl{D}}_\Psi^{w\tilde{w}}} =
\bmat{
    \mcl{A}-\mcl{B}\Lambda_c\mcl{D}_\Psi^{wz}\mcl{C} &
    -\mcl{B}\Lambda_c\mcl{C}_\Psi^w &
    \mcl{B}\Lambda_c}.
\end{align*}
For the second row,
\eqref{eqn:id2} and \eqref{eqn:id4} give
\begin{align*}
    \hat{\mcl{B}}_\Psi^z\Lambda_z
        &= \left(\mcl{B}_\Psi^z
            -\mcl{B}_\Psi^w(\mcl{D}_\Psi^w)^{-1}
             \mcl{D}_\Psi^{wz}\right)
            \Lambda_z\\
        &= \mcl{B}_\Psi^z\Lambda_z
            -\mcl{B}_\Psi^w(\mcl{D}_\Psi^w)^{-1}
             \mcl{D}_\Psi^{wz}\Lambda_z\\
        &= \mcl{B}_\Psi^z
            \left(I-\mcl{D}\Lambda_c\mcl{D}_\Psi^{wz}\right)
            -\mcl{B}_\Psi^w\Lambda_c\mcl{D}_\Psi^{wz}\\
        &= \mcl{B}_\Psi^z
            -\left(\mcl{B}_\Psi^z\mcl{D}+\mcl{B}_\Psi^w\right)
             \Lambda_c\mcl{D}_\Psi^{wz}.
\end{align*}
Its third entry follows from \eqref{eqn:id5}:
\begin{align*}
    &\hat{\mcl{B}}_\Psi^{\tilde{w}}
        +\hat{\mcl{B}}_\Psi^z\Lambda_z\mcl{D}
         \hat{\mcl{D}}_\Psi^{w\tilde{w}}\\
    &\quad =
        \mcl{B}_\Psi^w(\mcl{D}_\Psi^w)^{-1}
        +\left[
            \mcl{B}_\Psi^z
            -\left(\mcl{B}_\Psi^z\mcl{D}+\mcl{B}_\Psi^w\right)
             \Lambda_c\mcl{D}_\Psi^{wz}
         \right]\mcl{D}(\mcl{D}_\Psi^w)^{-1}\\
    &\quad =
        \left(\mcl{B}_\Psi^z\mcl{D}+\mcl{B}_\Psi^w\right)
        \left[
            (\mcl{D}_\Psi^w)^{-1}
            -\Lambda_c\mcl{D}_\Psi^{wz}\mcl{D}
             (\mcl{D}_\Psi^w)^{-1}
        \right]\\
    &\quad =
        \left(\mcl{B}_\Psi^z\mcl{D}+\mcl{B}_\Psi^w\right)\Lambda_c.
\end{align*}
The third row is obtained analogously:
\begin{equation*}
    \hat{\mcl{D}}_\Psi^z\Lambda_z
        = \mcl{D}_\Psi^z
           -\left(\mcl{D}_\Psi^z\mcl{D}+\mcl{D}_\Psi^{zw}\right)
            \Lambda_c\mcl{D}_\Psi^{wz} \qquad
    \hat{\mcl{D}}_\Psi^{z\tilde{w}}
        +\hat{\mcl{D}}_\Psi^z\Lambda_z\mcl{D}
         \hat{\mcl{D}}_\Psi^{w\tilde{w}}\\
    =
        \left(\mcl{D}_\Psi^z\mcl{D}+\mcl{D}_\Psi^{zw}\right)\Lambda_c.
\end{equation*}
Collecting these blocks and retaining $\tilde{w}$ as an output yields
\begin{equation}\label{eqn:congruence}
    \bmat{\mcl{T} \dot{\mbf{x}}(t) \\ \mcl{T}_\Psi\dot{\mbf{x}}_\Psi(t)\\ \tilde{z}(t) \\ \tilde{w}(t)} = 
    \bmat{\mcl{A} - \mcl{B}\Lambda_c\mcl{D}_\Psi^{wz}\mcl{C} & -\mcl{B}\Lambda_c\mcl{C}_\Psi^w & \mcl{B}\Lambda_c \\[4pt]
    \mcl{B}_\Psi^z \mcl{C} - (\mcl{B}_\Psi^z \mcl{D} + \mcl{B}_\Psi^w)\Lambda_c\mcl{D}_\Psi^{wz}\mcl{C} 
    & \mcl{A}_\Psi - (\mcl{B}_\Psi^z \mcl{D} + \mcl{B}_\Psi^w)\Lambda_c\mcl{C}_\Psi^w 
    & (\mcl{B}_\Psi^z \mcl{D} + \mcl{B}_\Psi^w)\Lambda_c \\[4pt]
    \mcl{D}_\Psi^z \mcl{C} - (\mcl{D}_\Psi^z \mcl{D} + \mcl{D}_\Psi^{zw})\Lambda_c\mcl{D}_\Psi^{wz}\mcl{C} 
    & \mcl{C}_\Psi^z - (\mcl{D}_\Psi^z \mcl{D} + \mcl{D}_\Psi^{zw})\Lambda_c\mcl{C}_\Psi^w 
    & (\mcl{D}_\Psi^z \mcl{D} + \mcl{D}_\Psi^{zw})\Lambda_c\\
    0&0&I}
    \bmat{\mbf{x}(t) \\ \mbf{x}_\Psi(t) \\ \tilde{w}(t)}
\end{equation}
The corresponding invertible congruence transformation is
\begin{equation}\label{eqn:congruencetf}
    \bmat{\mbf{x}(t) \\ \mbf{x}_\Psi(t) \\ \tilde{w}(t)} = 
    \bmat{I & 0 & 0\\ 
    0 & I & 0\\
    \mcl{D}_\Psi^{wz}\mcl{C} & \mcl{C}_\Psi^w & \Lambda_c^{-1}
    }\bmat{\mbf{x}(t) \\ \mbf{x}_\Psi(t) \\ {w}(t)}
    \qquad
    \bmat{\mbf{x}(t) \\ \mbf{x}_\Psi(t) \\ w(t)} = 
    \bmat{I & 0 & 0\\ 
    0 & I & 0\\
    -\Lambda_c\mcl{D}_\Psi^{wz} \mcl{C} & -\Lambda_c\mcl{C}_\Psi^w & \Lambda_c
    }\bmat{\mbf{x}(t) \\ \mbf{x}_\Psi(t) \\ \tilde{w}(t)}
\end{equation}
Applying \eqref{eqn:congruencetf} to \eqref{eqn:congruence} gives
\begin{equation}\label{eqn:congruence:transformed}
    \bmat{\mcl{T}\dot{\mbf{x}}(t)\\
          \mcl{T}_\Psi\dot{\mbf{x}}_\Psi(t)\\
          \tilde{z}(t)\\
          \tilde{w}(t)}
    =
    \bmat{
        \mcl{A} & 0 & \mcl{B}\\
        \mcl{B}_\Psi^z\mcl{C} & \mcl{A}_\Psi &
            \mcl{B}_\Psi^z\mcl{D}+\mcl{B}_\Psi^w\\
        \mcl{D}_\Psi^z\mcl{C} & \mcl{C}_\Psi^z &
            \mcl{D}_\Psi^z\mcl{D}+\mcl{D}_\Psi^{zw}\\
        \mcl{D}_\Psi^{wz}\mcl{C} & \mcl{C}_\Psi^w &
            \mcl{D}_\Psi^{wz}\mcl{D}+\mcl{D}_\Psi^w}
    \bmat{\mbf{x}(t)\\ \mbf{x}_\Psi(t)\\ w(t)}.
\end{equation}
The realization in \eqref{eqn:congruence:transformed} is the
direct augmentation of the \textbf{Primal System} with the \textbf{Primal Scaling}, and
hence is the \textbf{Scaled Primal System}. For necessity, applying \eqref{eqn:congruencetf} to
\eqref{eqn:congruence:transformed} recovers \eqref{eqn:congruence},
and reversing \eqref{eqn:id1}--\eqref{eqn:id5} recovers
\eqref{eqn:potapov}. The dual case follows by applying the same
argument to the \textbf{Dual System} and the \textbf{Dual Scaling}.
\end{proof}
\begin{theorem}[Scaled IO duality]\label{thm:dualscaledIO}
	Assume the following conditions hold
	\begin{itemize}
		\item The \textbf{Primal I/O Scaling} is a PI operator and invertible.
		\item $\mcl{S}_{11} - \mcl{D}\mcl{S}_{21}$ and $\bar{\mcl{S}}_{11} - \mcl{D}^*\bar{\mcl{S}}_{21}$ are invertible.
	\end{itemize}
	The following statements are equivalent.
	\begin{enumerate}
		\item For all $\tilde{w}\in\RL^{\mbf{n}_w}$, the \textbf{Scaled Primal System}
		      satisfies $\|P_\tau\tilde{z}\| \leq \|P_\tau\tilde{w}\|$.
		\item For all $\tilde{\bar{w}}\in\RL^{\mbf{n}_z}$, the \textbf{Scaled Dual System}
		      satisfies $\|P_\tau\tilde{\bar{z}}\| \leq \|P_\tau\tilde{\bar{w}}\|$.
	\end{enumerate}
\end{theorem}
\begin{proof}
	See Appendix \ref{proof:theorem2}.
\end{proof}
